\section{Convergence Analysis}
\label{sec:convergence}

\subsection{Formal Convergence Theorem}

We state the convergence result in the language of our application.

\begin{definition}[Coarse Acceptance Function]
  \label{def:accept-fn}
  For $\alpha \in [0.05, 0.95]$, let
  $A(\alpha) = \mathrm{E}_\alpha[\text{coarse acceptance rate}]$
  be the expected fraction of coarse-move steps that are accepted (projection
  succeeds and rebalancer converges within 200 iterations), under the
  stationary distribution of the \ams{} chain with fixed $\alpha$.
  Let $\tau = 0.30$ be the target.
\end{definition}

\begin{definition}[Optimal Alpha]
  \label{def:alpha-star}
  The \emph{optimal alpha} is $\alpha^* \in [0.05, 0.95]$ such that
  $A(\alpha^*) = \tau$, if such a value exists.
\end{definition}

\begin{theorem}[Convergence of \ams{}, informal]
  \label{thm:convergence}
  Assume:
  \begin{enumerate}
    \item[\textup{(M)}] \emph{Monotonicity}: $A(\alpha)$ is non-decreasing in $\alpha$ on
      $[0.05, 0.95]$ and there exists $\alpha^* \in (0.05, 0.95)$ with
      $A(\alpha^*) = \tau$.
    \item[\textup{(B)}] \emph{Bounded noise}: the variance of the observed coarse
      acceptance rate in any window of $\Delta$ steps is bounded uniformly in
      $\alpha$ and in the chain history.
  \end{enumerate}
  Then the \ams{} alpha trace $\{\alpha_t\}$ converges to $\alpha^*$ in
  mean square as $t \to \infty$.
\end{theorem}

\begin{proof}[Proof sketch]
  At each adaptation round $t$, the update is
  \[
    \alpha_{t+1} = \mathrm{clip}\!\bigl(\alpha_t + \gamma_t(\hat{r}_t - \tau),\; 0.05,\; 0.95\bigr),
  \]
  where $\hat{r}_t = r_t + \varepsilon_t$ is the observed acceptance rate,
  $r_t = A(\alpha_t)$ is the expected rate, and $\varepsilon_t$ is zero-mean
  noise with bounded variance (assumption B).
  Writing $h(\alpha) = A(\alpha) - \tau$, the update (ignoring the clip) is:
  \[
    \alpha_{t+1} = \alpha_t + \gamma_t h(\alpha_t) + \gamma_t \varepsilon_t.
  \]
  By assumption M, $h(\alpha)$ is non-decreasing with unique root $\alpha^*$.
  The step-size sequence $\gamma_t = \gamma_0/t^{3/4}$ satisfies
  $\sum \gamma_t = \infty$ and $\sum \gamma_t^2 < \infty$
  (Section~\ref{sec:background}).
  By the Robbins-Monro theorem~\citep{robbins1951}, the unclipped iterate
  converges in mean square to $\alpha^*$.
  The clip $[0.05, 0.95]$ contains $\alpha^*$ by assumption M
  (the root exists in the interior of the clip range); the clip therefore
  activates only finitely often almost surely, and convergence is unaffected.
\end{proof}

\paragraph{Empirical verification of monotonicity.}
We verify the monotonicity assumption empirically: running MultiScale with
fixed $\alpha \in \{0.10, 0.50\}$ on NC for 500 steps each, we observe
$p_{\text{rebalance}}(0.10) = 0.94$ and $p_{\text{rebalance}}(0.50) = 0.89$
--- a 5-point variation, consistent with the assumption that $A(\alpha)$ is
approximately monotone over the operational range $[0.05, 0.95]$.
While this two-point check does not rule out non-monotonicity at
intermediate values, it establishes the direction and approximate magnitude
of the relationship.

\paragraph{Caveat on the monotonicity assumption.}
Assumption~(M) requires that the expected coarse acceptance rate is
non-decreasing in $\alpha$.
This is intuitively plausible: a higher $\alpha$ proposes more coarse moves,
but the \emph{acceptance rate per coarse proposal} does not depend on $\alpha$
(it depends on the rebalancer's success probability, which is a function of
the plan, not of $\alpha$).
Thus $A(\alpha) = \alpha \cdot A_0 / \alpha = A_0$ is actually constant in
$\alpha$ under the stationary assumption --- which means the root
$A(\alpha^*) = \tau$ either exists everywhere or nowhere.

This argument exposes a subtlety: $A(\alpha)$ as defined is the
\emph{fraction of all steps} (including fine steps) that are accepted coarse
moves.
This equals $\alpha \cdot p_{\text{rebalance}}$, where $p_{\text{rebalance}}$
is the rebalancer success probability.
Then $A(\alpha) = \tau$ implies $\alpha^* = \tau / p_{\text{rebalance}}$.
This is indeed monotone in $\alpha$ (actually linear), and the Robbins-Monro
update converges to $\alpha^* = \tau / p_{\text{rebalance}}$ which achieves
a coarse-acceptance fraction of $\tau$.

However, $p_{\text{rebalance}}$ itself may depend weakly on $\alpha$ (at
higher $\alpha$, more coarse moves are proposed, and the plan state seen by
the rebalancer is influenced by the history of accepted coarse moves).
This coupling is empirically small but is not formally characterised.
We treat this as an empirically verified assumption and not a formal guarantee.

\subsection{Empirical Alpha Traces}

Table~\ref{tab:alpha-convergence} and the narrative below summarise the
empirical alpha trace for North Carolina, Texas, and Iowa, each starting at
$\alpha_0 = 0.30$, with $T = 5{,}000$ steps, $\Delta = 50$, $\gamma_0 = 0.10$.

\begin{table}[h]
  \centering
  \caption{Empirical alpha convergence summary.}
  \label{tab:alpha-convergence}
  \begin{tabular}{lrrrrr}
    \toprule
    State & $k$ & Counties & Tracts & $\alpha_0$ & $\alpha^*$ (empirical) \\
    \midrule
    NC & 14 & 100 & 2{,}195 & 0.30 & $\approx 0.28$ \\
    TX & 38 & 254 & 5{,}265 & 0.30 & $\approx 0.33$ \\
    IA &  4 &  99 &   825 & 0.30 & $\approx 0.17$ \\
    \bottomrule
  \end{tabular}
\end{table}

\paragraph{North Carolina ($k = 14$).}
NC has a county/tract ratio of $2{,}195 / 100 \approx 22.0$, and $k = 14$
districts.
The county structure is reasonably uniform (no single county dominates);
the rebalancer succeeds at a moderate rate.
The alpha trace starts at 0.30, oscillates in the first 20 adaptation rounds
as the chain warms up, and stabilises near 0.28 --- slightly below the target,
reflecting a small bias from the discrete clip.

\paragraph{Texas ($k = 38$).}
TX has the highest $k$ and the most counties.
The county/tract ratio ($\approx 20.7$) is similar to NC, but the larger $k$
means the rebalancer has more districts to balance.
Despite this, the coarse acceptance rate with $\alpha = 0.30$ is slightly
above target because TX's flat West Texas geography creates counties that
are easy to project.
The alpha trace converges to $\approx 0.33$, reflecting that the chain can
tolerate slightly more coarse moves than the default target.

\paragraph{Iowa ($k = 4$).}
IA has the smallest $k$ and a relatively small county/tract ratio ($\approx 8.3$).
With only four districts, the rebalancer must achieve a tight population
balance on county-level proposals with coarse counties.
The acceptance rate at $\alpha = 0.30$ is well below 0.30 (many coarse
proposals fail rebalancing), so the alpha trace decreases monotonically to
$\approx 0.17$.
Convergence is fastest for IA because the feedback signal is strong (many
rejections drive $\alpha$ down decisively).

\paragraph{Convergence speed.}
In all three states, $|\alpha_t - \alpha^*| < 0.05$ after approximately
40--60 adaptation rounds (2{,}000--3{,}000 steps).
The remaining steps run with a near-optimal $\alpha$, achieving mixing
comparable to or better than fixed-$\alpha$ \ms{} with the state-specific
optimal $\alpha$.

\subsection{Convergence Diagnostics in the Manifest}

The run manifest records \texttt{final\_alpha} and issues a warning if
convergence appears incomplete:
\begin{verbatim}
NOTE: alpha has not converged (|final - target| = 0.183);
consider more steps or a larger gamma_0.
\end{verbatim}
This warning fires when $|\texttt{final\_alpha} - \tau| > 0.15$ after at least
$T \ge 500$ chain steps (i.e., at least $\lfloor 500/\Delta \rfloor = 10$
adaptation rounds with default $\Delta = 50$).
It is a diagnostic, not a hard error: the plan is still valid and auditable.
